\title{DeepSeekMath: Pushing the Limits of Mathematical Reasoning in Open Language Models}

\begin{abstract}
The intricacies and structured aspects of mathematical reasoning present notable hurdles for language models. Here, we present DeepSeekMath~7B, which extends the pre-training of DeepSeek-Coder-Base-v1.5 7B utilizing 120B math-specific tokens gathered from Common Crawl, as well as data in both natural language and code. Without dependence on auxiliary toolkits or ensemble voting, DeepSeekMath~7B attains a strong 51.7\% on the competition-level MATH benchmark, nearing the performance achieved by Gemini-Ultra and GPT-4. When employing self-consistency over 64 samples, DeepSeekMath~7B further improves to 60.9\% accuracy on MATH. The remarkable mathematical reasoning proficiency of DeepSeekMath~is largely due to two principal contributors: first, leveraging the considerable wealth of public web data via a rigorously designed data selection pipeline; and second, introducing Group Relative Policy Optimization (GRPO)—a modified form of Proximal Policy Optimization (PPO)—which both advances mathematical reasoning capacity and effectively reduces PPO's memory footprint.
\end{abstract}

\section{Introduction}

Large language models (LLMs) have transformed methodologies for mathematical reasoning within artificial intelligence, leading to notable progress on tasks such as the quantitative reasoning benchmark \citep{MATH} and the geometry reasoning benchmark \citep{trinh2024solving}. Additionally, these systems have played a key role in enabling humans to tackle challenging mathematical problems \citep{tao}. Despite these developments, state-of-the-art systems like GPT-4 \citep{gpt4} and Gemini-Ultra \citep{gemini} remain closed to the public, while those open-source models that are available demonstrate substantially inferior performance.

This paper presents DeepSeekMath, a specialized language model tailored for mathematical reasoning, which delivers substantially stronger performance than current open-source systems and approaches the proficiency of GPT-4 across academic test suites. Central to this achievement is the construction of the DeepSeekMath Corpus, an extensive, high-quality pre-training dataset totaling 120B mathematical tokens. Compilation of this dataset involves leveraging the Common Crawl (CC) and applying a fastText-driven classifier \citep{joulin2016fasttext}. In its initial phase, the classifier is trained on positive samples from OpenWebMath \citep{openwebmath} and negatives drawn from varied web sources. Thereafter, the model utilizes this classifier to identify more potential positives within CC, with these candidates being refined further through manual annotation. The improved dataset feeds back into subsequent iterations of the classifier, thereby enhancing its precision. Empirical validation shows the resulting corpus is of considerable quality: the DeepSeekMath-Base 7B model attains 64.2\% accuracy on GSM8K \citep{gsm8k} and 36.2\% on the challenging MATH competition set \citep{MATH}, outperforming the Minerva 540B baseline \citep{minerva}. Owing to the corpus’s multilingual composition, we also observe marked gains on Chinese mathematical evaluation benchmarks \citep{wei2023cmath,agieval}. We consider our efforts in mathematical corpus development to lay important groundwork for subsequent advances, acknowledging ample prospects for future refinement.

DeepSeekMath-Base is built upon DeepSeek-Coder-Base-v1.5 7B \citep{deepseek-coder}, since our findings suggest that leveraging a code-centric pre-trained model offers advantages over utilizing a generic LLM foundation. In addition, we find that training on mathematical tasks yields improvements not only in the model's performance on mathematical benchmarks, but also enhances results on general reasoning tests such as MMLU \citep{mmlu} and BBH \citep{bbh}. This suggests that the benefits of math-oriented training extend to broader reasoning capabilities as well.

Following the pre-training phase, we perform mathematical instruction fine-tuning on DeepSeekMath-Base utilizing datasets incorporating chain-of-thought \citep{cot}, program-of-thought \citep{pot,pal}, and tool-integrated reasoning \citep{tora} approaches. The final model, DeepSeekMath-Instruct 7B, surpasses all other 7B models and achieves performance on par with instruction-tuned open-source models at the 70B scale.

In addition, we propose Group Relative Policy Optimization (GRPO), an alternative reinforcement learning (RL) algorithm inspired by Proximal Policy Optimization (PPO) \citep{schulman2017proximal}. GRPO eliminates the need for a critic model, instead deriving the baseline directly from aggregated group scores, thereby substantially reducing the computational demands during training. Leveraging only a limited portion of English instruction-tuning data, GRPO delivers notable performance gains over the already competitive DeepSeekMath-Instruct, with improvements observed in both in-domain tasks (GSM8K: 82.9\% $\rightarrow$ 88.2\%, MATH: 46.8\% $\rightarrow$ 51.7\%) and out-of-domain challenges (for instance, CMATH: 84.6\% $\rightarrow$ 88.8\%) during the RL stage. Furthermore, we present a coherent framework for interpreting existing approaches—including Rejection Sampling Fine-Tuning (RFT) \citep{yuan2023scaling}, Direct Preference Optimization (DPO) \citep{dpo}, PPO, and GRPO—demonstrating that these techniques fundamentally correspond to either direct or streamlined forms of RL. Through comprehensive experimental analysis—comparing, for example, online and offline learning, outcome versus process-based supervision, and single-turn against iterative RL—we thoroughly explore the core components of this overarching framework. Finally, we elucidate the mechanisms through which our RL strategy enhances instruction-tuned model performance, and we discuss prospective research directions for more efficient RL development grounded in this unified perspective.

\subsection{Contributions}
We present a scalable approach to math pre-training and conduct a thorough investigation and assessment of reinforcement learning methodologies.

\textbf{Math Pre-Training at Scale}

\begin{itemize}[topsep=0pt]
    \item This study demonstrates that the openly available Common Crawl dataset serves as a rich resource for mathematical applications. Through a carefully developed data filtering process, we assemble the DeepSeekMath Corpus, comprising 120B tokens sourced from web pages explicitly selected for mathematical relevance. This corpus is nearly sevenfold larger than the math-focused web pages utilized in Minerva \citep{minerva} and exceeds the size of the recently published OpenWebMath \citep{openwebmath} by a factor of nine.

\item DeepSeekMath-Base 7B, our pre-trained base model, demonstrates similar results to Minerva 540B \citep{minerva}, suggesting that parameter count alone does not solely determine mathematical reasoning proficiency. Strong outcomes can also be attained by more compact models when pre-trained on high-quality datasets.

\item We present results from our experiments involving math training. Pretraining models on code before exposing them to math tasks enhances their performance in solving mathematical problems, regardless of whether they utilize external tools. This provides some insight into the ongoing debate: \textit{does code training enhance reasoning skills?} Based on our observations, we find that it does—specifically in the realm of mathematical reasoning.

\item While it is a standard practice to train models using arXiv papers, particularly within the mathematical domain, this strategy does not yield any significant enhancement across the mathematical benchmarks utilized in this study.
\end{itemize}

\textbf{Exploration and Analysis of Reinforcement Learning}

\begin{itemize}[topsep=0pt]
    \item We present Group Relative Policy Optimization (GRPO), a reinforcement learning approach that operates without a critic by leveraging group-based score baselines. This design considerably lowers the computational demands required for training when compared to Proximal Policy Optimization (PPO), while maintaining efficiency and effectiveness.
    \item Through empirical evaluation, we show that applying GRPO substantially improves the capabilities of our instruction-tuned model, DeepSeekMath-Instruct, relying solely on instruction-tuning data. Additionally, the method yields better generalization to out-of-domain scenarios throughout reinforcement learning.
    \item We offer a comprehensive framework that places various methods—such as RFT, DPO, PPO, and GRPO—within a unified conceptual context. Moreover, we perform a broad range of experiments, examining facets such as online versus offline optimization, outcome- versus process-based supervision, and single-turn versus iterative reinforcement learning, aiming to thoroughly analyze the critical aspects underlying this paradigm.
    \item Grounded in our unified framework, we investigate the underlying causes for reinforcement learning’s effectiveness and outline several promising avenues for achieving more efficient reinforcement learning in large language models.
\end{itemize}

\subsection{Summary of Evaluations and Metrics}

\begin{itemize}[topsep=0pt]
    \item \textbf{English and Chinese Mathematical Reasoning}: Our study presents a thorough evaluation of model performance using a variety of English and Chinese mathematical benchmarks, encompassing tasks from elementary to collegiate complexity. The English benchmarks utilized are GSM8K \citep{gsm8k}, MATH \citep{MATH}, SAT \citep{llemma}, OCW Courses \citep{minerva}, and MMLU-STEM \citep{mmlu}, while the Chinese datasets comprise MGSM-zh \citep{mgsm}, CMATH \citep{wei2023cmath}, Gaokao-MathCloze \citep{agieval}, and Gaokao-MathQA \citep{agieval}. We assess the models’ proficiency both in composing fully self-explanatory text-based solutions without external assistance and in problem-solving via Python programming.

On evaluations using English benchmarks, DeepSeekMath-Base achieves results comparable to the proprietary Minerva 540B \citep{minerva}, and outperforms all existing open-source base models—including Mistral 7B \citep{mistral} and Llemma-34B \citep{llemma}—regardless of whether these models have been specifically pre-trained on mathematics-related data, frequently by a substantial margin. Interestingly, DeepSeekMath-Base attains even greater performance on Chinese benchmarks, which can be attributed to our departure from prior approaches \citep{minerva,llemma} that focus solely on English mathematical pre-training corpora; our method incorporates high-quality mathematical data from additional languages. Through instruction tuning tailored to mathematics, along with reinforcement learning strategies, the derived DeepSeekMath-Instruct and DeepSeekMath-RL models set a new standard in the open-source arena by achieving over 50\% accuracy for the first time on the competition-level MATH dataset.

\item \textbf{Formal Mathematics}: DeepSeekMath-Base is assessed on its ability to perform informal-to-formal theorem proving, following the methodology introduced by \citep{dsp_proof}. Evaluations are carried out on the miniF2F benchmark \citep{minif2f}, utilizing Isabelle \citep{isabelle} as the proof assistant. The results indicate that DeepSeekMath-Base achieves impressive performance in few-shot autoformalization scenarios.

\item \textbf{Natural Language Understanding, Reasoning, and Code}: To thoroughly assess the model’s proficiency in general language comprehension, complex reasoning, and programming, we benchmark DeepSeekMath-Base on several diverse tasks. This includes the Massive Multitask Language Understanding (MMLU) suite \citep{mmlu}, composed of 57 distinct multiple-choice topics, BIG-Bench Hard (BBH) \citep{bbh}, a collection of 23 tasks mainly targeting multi-step reasoning, alongside HumanEval \citep{codex} and MBPP \citep{mbpp}, both of which are established benchmarks for evaluating code generation abilities. Our findings suggest that pre-training on mathematical data leads to improvements not only in reasoning but also in general language understanding.

\section{Math Pre-Training}

\subsection{Data Collection and Decontamination} This section describes the methodology used to assemble the DeepSeekMath Corpus leveraging Common Crawl data. As illustrated in Figure \ref{fig:data_collect}, we introduce an iterative workflow for the systematic extraction of a large-scale mathematics corpus from Common Crawl, which initially begins with a seed corpus comprising a limited, yet high-quality, math-centric dataset. It should be emphasized that this methodology can also be extended to domains beyond mathematics, including programming.

Initially, we utilize OpenWebMath \citep{openwebmath}, a well-curated repository of mathematical content from the web, as our primary seed dataset. From this foundation, we construct a fastText model \citep{joulin2016fasttext} aimed at retrieving additional web pages with characteristics similar to those in OpenWebMath. For the training procedure, we draw a random sample of 500,000 entries from the seed as positive instances, complemented by 500,000 web documents taken from Common Crawl to serve as negative samples. Model training leverages a publicly available library\footnote{\url{https://fasttext.cc}}, with hyperparameters set as follows: vector dimensionality at 256, a learning rate of 0.1, a word n-gram maximum length of 3, a minimum threshold of 3 word occurrences, and 3 epochs for model convergence. The initial bulk from Common Crawl undergoes reduction via both URL deduplication and near-duplicate filtering, yielding a set of 40 billion HTML pages. Using the fastText model, we then identify and extract mathematical content from this deduplicated data. To refine the results and eliminate pages with subpar mathematical material, the web pages are ranked based on the model's predicted scores, after which only the highest-ranked documents are retained. The scale of retained data is measured through pre-training trials considering the largest 40B, 80B, 120B, and 160B token subsets. For this initial stage, the selection is restricted to the top 40B tokens.

Following the initial round of data gathering, a substantial portion of mathematical web pages remains uncaptured, primarily due to insufficient diversity in the positive samples used to train the fastText model. To address this, we broaden the seed corpus by incorporating more mathematical web resources, enhancing the model's capacity. Our method involves partitioning the entire Common Crawl dataset into distinct domains, where each domain comprises web pages sharing an identical base URL. For every domain, we determine the proportion of pages retrieved during the first iteration. Domains with over 10\% of pages retrieved are designated as math-related (for instance, \url{mathoverflow.net}). We then manually curate and label the specific URLs tied to mathematical content within these domains (such as \url{mathoverflow.net/questions}). Any previously uncollected web pages linked to these URLs are subsequently integrated into the seed corpus. This process increases the pool of positive samples, facilitating the training of a more effective fastText model that can recognize a greater volume of mathematical content in later cycles. Upon completing four data collection rounds, we assemble a dataset containing 35.5M mathematical web pages, amounting to 120B tokens. By the fourth iteration, we observe that approximately 98\% of the content had already been retrieved in the prior iteration, leading us to terminate further data collection.

In order to prevent benchmark leakage, we adhere to the methodology outlined in \cite{deepseek-coder}, excluding web pages that feature questions or answers from established English mathematical benchmarks such as GSM8K \citep{gsm8k} and MATH \citep{MATH}, as well as Chinese benchmarks including CMATH \citep{wei2023cmath} and AGIEval \citep{agieval}. The process for filtering is as follows: any segment of text within our math training corpus that contains a 10-gram sequence perfectly matching any substring from these evaluation benchmarks is discarded. Additionally, for benchmark content consisting of fewer than 10 but at least 3 grams, we apply exact matching to similarly remove any potentially contaminated web pages.

\subsection{Validating the Quality of the DeepSeekMath~Corpus}

We run pre-training experiments to investigate how the DeepSeekMath~Corpus is compared with the recently released math-training corpora:

\begin{itemize}[topsep=0pt]
    \item \textbf{MathPile} \citep{mathpile}: This is a composite dataset comprising 8.9B tokens, drawn from multiple sources such as textbooks, Wikipedia, ProofWiki, CommonCrawl, StackExchange, and arXiv. Notably, arXiv makes up over 85\% of the total token count;
    \item \textbf{OpenWebMath} \citep{openwebmath}: Constructed by selecting mathematical materials from CommonCrawl, this corpus consists of 13.6B tokens;
    \item \textbf{Proof-Pile-2} \citep{llemma}: This collection merges OpenWebMath, AlgebraicStack (containing 10.3B tokens of math code), and arXiv papers (totaling 28.0B tokens). For experimental purposes with Proof-Pile-2, we adhere to the arXiv:Web:Code ratio of 2:4:1 as prescribed by \cite{llemma}.
\end{itemize}

\subsubsection{Training Setting}
\label{sec:quality-policy}
We perform math-centric fine-tuning on a foundational language model comprised of 1.3B parameters, architecturally aligned with the DeepSeek LLMs \citep{deepseek-llm}, and referred to as DeepSeek-LLM 1.3B. Each mathematical dataset is used to independently train a separate model for a total of 150B tokens. All training runs leverage the resource-efficient and streamlined HAI-LLM \citep{haillm} system. Adhering to the protocol established by DeepSeek LLMs, we utilize the AdamW optimizer \citep{adamW}, assigning $\beta_1=0.9$, $\beta_2=0.95$, and $\mathrm{weight\_decay}=0.1$. Our learning rate scheduling follows a multi-stage scheme: the rate is linearly ramped up during the first 2,000 warmup steps, decays to 31.6\% of its peak value after reaching 80\% completion of total training, and is further reduced to 10.0\% of the maximum by the 90\% mark. The highest learning rate allowed is set at 5.3e-4. We conduct training using batch sizes of 4 million tokens and employ a 4K context window.

\subsubsection{Evaluation Results}

\textbf{The DeepSeekMath~Corpus is of high quality, covers multilingual mathematical content, and is the largest in size.}

\begin{itemize}[topsep=0pt]
    \item \textbf{High-quality}: 
    We assess downstream capabilities across 8 mathematics-focused benchmarks utilizing few-shot chain-of-thought prompting, as described in \cite{cot}. The data in Table \ref{tab:corpora_comparison} clearly reveal superior performance by models trained on the DeepSeekMath~Corpus. Additionally, Figure \ref{fig:corpora_comparisons} highlights that, at 50B tokens (equivalent to a complete pass through Proof-Pile-2), models utilizing DeepSeekMath substantially surpass those based on Proof-Pile-2, signifying a higher overall quality in the DeepSeekMath~Corpus.
    \item \textbf{Multilingual}:
    The DeepSeekMath~Corpus contains multilingual content, with English and Chinese comprising the majority of entries. According to Table \ref{tab:corpora_comparison}, model training with DeepSeekMath~Corpus results in significant gains in mathematical reasoning tasks for both English and Chinese. In contrast, most available mathematics corpora, which focus mainly on English, provide minimal benefit and, in some cases, even degrade Chinese language reasoning performance.
    \item \textbf{Large-scale}:
    DeepSeekMath~Corpus greatly exceeds the size of prior mathematical datasets. As seen in Figure \ref{fig:corpora_comparisons}, DeepSeek-LLM 1.3B, when trained on this corpus, achieves a much faster learning trajectory and demonstrates more sustained advancements. Meanwhile, the comparative baseline corpora are substantially smaller, with their examples recycled numerous times during training, resulting in early stagnation of model accuracy.
\end{itemize}

\subsection{Training and Evaluating DeepSeekMath-Base 7B}

This section presents DeepSeekMath-Base 7B, a foundational model distinguished by its robust mathematical reasoning capabilities. The model is initialized with DeepSeek-Coder-Base-v1.5 7B \citep{deepseek-coder} and subsequently trained over 500B tokens. The data composition is as follows: 56\% originates from the DeepSeekMath Corpus, 4\% from AlgebraicStack, 10\% from arXiv, 20\% from Github code repositories, and the final 10\% comprises natural language materials sourced from Common Crawl in both English and Chinese. The primary training configuration adheres to the guidelines laid out in Section \ref{sec:quality-policy}, with two exceptions: the learning rate peak is adjusted to 4.2e-4, and a batch size of 10M tokens is utilized.

We present an in-depth evaluation of DeepSeekMath-Base 7B’s mathematical proficiencies, examining its capacity to independently generate complete mathematical solutions, utilize tools for problem-solving, and execute formal theorem proving. In addition to its mathematical aptitude, we deliver a broader overview of the base model, highlighting its effectiveness in natural language understanding, reasoning tasks, and programming abilities.

\paragraph{Mathematical Problem Solving with Step-by-Step Reasoning}

We assess the ability of DeepSeekMath-Base to solve mathematical problems through few-shot chain-of-thought prompting \citep{cot}, utilizing eight different benchmarks in both English and Chinese. These benchmarks include tasks requiring quantitative reasoning, such as GSM8K \citep{gsm8k}, MATH \citep{MATH}, and CMATH \citep{wei2023cmath}, as well as multiple-choice assessments like MMLU-STEM \citep{mmlu} and Gaokao-MathQA \citep{agieval}. Collectively, they span a wide range of mathematical disciplines, from elementary-level to college-level problems.

Table \ref{tab:base_math_cot} demonstrates that DeepSeekMath-Base 7B consistently achieves top results on all eight evaluated benchmarks when compared with other open-source base models, such as the popular general-purpose Mistral 7B \citep{mistral} and the newly introduced Llemma 34B \citep{llemma}, the latter of which is additionally trained on the Proof-Pile-2 dataset \citep{llemma}. Particularly striking is DeepSeekMath-Base’s performance on the MATH dataset, where it exceeds all existing open-source base models by a margin greater than 10\% absolute, and even surpasses Minerva 540B \citep{minerva}—a proprietary model that is 77 times larger, built upon PaLM \citep{palm}, and further specialized on mathematical literature.

\paragraph{Mathematical Problem Solving with Tool Use} We assess the ability of models to solve mathematical problems using external tools on the GSM8K and MATH benchmarks, employing few-shot program-of-thought prompting as outlined by \citep{pot,pal}. In this setup, models address each question by generating a Python script, leveraging libraries like \textit{math} and \textit{sympy} to handle complex calculations. The output generated by executing this script is taken as the predicted answer. According to the results in Table \ref{tab:base_math_pot_proof}, DeepSeekMath-Base 7B surpasses the previous top-performing model, Llemma 34B.

\paragraph{Formal Mathematics}

Automating the verification of formal proofs offers significant advantages in promoting both the correctness and dependability of mathematical argumentation, while also improving workflow efficiency—a focus that has seen heightened interest in the literature of late. In our assessment, we use DeepSeekMath-Base 7B to address the informal-to-formal proof task introduced by \citep{dsp_proof}, which entails generating a complete formal proof given an informal problem description, its formal equivalent, and a corresponding informal proof sketch. For this purpose, the evaluation is conducted using the miniF2F \citep{minif2f} benchmark, a suite specifically curated for formal reasoning on Olympiad-style mathematics problems, where for each instance we prompt the model few-shot to produce an Isabelle-formalized proof. Adopting the methodology from \cite{dsp_proof}, the approach combines model-generated proof outlines with the automated prover Sledgehammer \citep{sledgehammer}, which systematically fills in missing inferential steps. The empirical results summarized in Table \ref{tab:base_math_pot_proof} indicate that DeepSeekMath-Base 7B achieves notably robust outcomes in automatic proof formalization.

\paragraph{Natural Language Understanding, Reasoning, and Code}

We assess natural language understanding using the MMLU benchmark \citep{mmlu}, reasoning ability via BBH \citep{bbh}, and coding proficiency with HumanEval \citep{codex} and MBPP \citep{mbpp}. Table \ref{tab:understanding_reasoning_code} presents the results, where DeepSeekMath-Base 7B demonstrates marked improvements over its predecessor, DeepSeek-Coder-Base-v1.5 \citep{deepseek-coder}, on both MMLU and BBH. These findings highlight the beneficial effects of mathematics-focused training for both understanding and reasoning in language models. Furthermore, by incorporating code tokens during continual training, DeepSeekMath-Base 7B successfully preserves the coding benchmark scores achieved by DeepSeek-Coder-Base-v1.5. Taken together, DeepSeekMath-Base 7B achieves notably higher scores than the widely-used Mistral 7B \citep{mistral} on the reasoning and coding tasks examined.

\section{Supervised Fine-Tuning}

\subsection{SFT Data Curation}

We develop a mathematical instruction-tuning dataset that includes problems in both English and Chinese, spanning multiple mathematical domains and a range of difficulty levels. Each problem is accompanied by a corresponding solution, formatted in chain-of-thought (CoT) \citep{cot}, program-of-thought (PoT) \citep{pot,pal}, or tool-integrated reasoning \citep{tora} styles. The dataset comprises a total of 776K training samples.

\begin{itemize}[topsep=0pt]
    \item \textbf{English mathematical datasets}: We provide tool-augmented solution annotations for problems from GSM8K and MATH, and further include select items from MathInstruct \citep{MathInstruct} as well as the Lila-OOD training set \citep{lila}, where answers are derived through either CoT or PoT methods. Our assembled English datasets encompass a broad spectrum of mathematical areas, including algebra, probability, number theory, calculus, and geometry.
    \item \textbf{Chinese mathematical datasets}: Our Chinese dataset comprises K-12 mathematics questions from 76 distinct sub-domains, such as linear equations, each accompanied by solutions formulated using both CoT reasoning and tool-integrated approaches.
\end{itemize}

\subsection{Training and Evaluating DeepSeekMath-Instruct 7B}

In this section, we present DeepSeekMath-Instruct 7B, which is derived from DeepSeekMath-Base through specialized mathematical instruction tuning. For training, examples are concatenated in a random order to fill a context window of up to 4K tokens. The model is optimized for 500 steps, utilizing a batch size of 256 and maintaining a fixed learning rate of 5e-5 throughout the process.

We evaluate models' mathematical performance both without and with tool use, on 4 quantitative reasoning benchmarks in English and Chinese. We benchmark our model against the leading models of the time:

\begin{itemize}[topsep=0pt]
    \item \textbf{Closed-source models} comprise: (1) the GPT lineup, with GPT-4 \citep{gpt4} and the GPT-4 Code Interpreter~\footnote{\url{https://openai.com/blog/chatgpt-plugins\#\#code-interpreter}} representing the most advanced options, (2) Gemini Ultra and Pro \citep{gemini}, (3) Inflection-2 \citep{inflection-2}, (4) Grok-1~\footnote{\url{https://x.ai/model-card}},
    in addition to newer offerings by Chinese developers, including (5) Baichuan-3~\footnote{\url{https://www.baichuan-ai.com}},
    and (6) the most recent version of GLM, GLM-4~\footnote{\url{https://open.bigmodel.cn/dev/api\#glm-4}}

    from the GLM family \citep{glm}.

These models are intended for broad applications, and most have been subjected to various alignment techniques.  
\item \textbf{Open-source models} encompass: general models such as (1) DeepSeek-LLM-Chat 67B \citep{deepseek-llm}, (2) Qwen 72B \citep{qwen}, (3) SeaLLM-v2 7B \citep{seallm}, and (4) ChatGLM3 6B \citep{chatglm3}. There are also models specifically improved for mathematical tasks, including (5) InternLM2-Math 20B~\footnote{\url{https://github.com/InternLM/InternLM-Math}}, developed by further training InternLM2 on mathematical content followed by instruction tuning, (6) Math-Shepherd-Mistral 7B, which utilizes PPO-based training \citep{schulman2017proximal} on top of Mistral 7B \citep{mistral} and incorporates a process-level reward mechanism, (7) the WizardMath family \citep{wizardmath}, which enhances the mathematical abilities of Mistral 7B and Llama-2 70B \citep{llama2} through evolve-instruct (a variant of instruction tuning that uses AI-generated instructions) and PPO training, drawing most of its training problems from GSM8K and MATH, (8) MetaMath 70B \citep{metamath}, produced by fine-tuning Llama-2 70B using an expanded dataset based on GSM8K and MATH, (9) ToRA 34B \cite{tora}, a version of CodeLlama 34B optimized for integrated mathematical tool use, and (10) MAmmoTH 70B \citep{MathInstruct}, which is Llama-2 70B refined through instruction tuning utilizing the MathInstruct collection.

Table \ref{tab:sft_rl_math} presents results under evaluation conditions that prohibit tool use, where DeepSeekMath-Instruct 7B displays robust capabilities in stepwise reasoning. In particular, on the MATH dataset—characterized by competition-level problems—our model achieves an absolute improvement of at least 9\% over all other open-source models and the majority of closed-source systems (such as Inflection-2 and Gemini Pro). This superiority holds even in comparison to models with significantly larger parameter counts (like Qwen 72B) or those explicitly optimized through reinforcement learning techniques targeting mathematical reasoning (for example, WizardMath-v1.1 7B). Although DeepSeekMath-Instruct performs comparably to leading proprietary Chinese models such as GLM-4 and Baichuan-3 on the MATH dataset, it does not yet match the results attained by GPT-4 or Gemini Ultra.

In the assessment scenario permitting models to employ both natural language reasoning and programmatic tool usage to address problems, DeepSeekMath-Instruct 7B achieves nearly 60\% accuracy on MATH, outperforming every other open-source model available. Across additional benchmarks, our model demonstrates competitive results when compared to DeepSeek-LLM-Chat 67B, the previous state-of-the-art, despite being one-tenth the size.

\subsection{Group Relative Policy Optimization} After the Supervised Fine-Tuning (SFT) phase, reinforcement learning (RL) has demonstrated significant capability in enhancing the mathematical reasoning performance of LLMs \citep{wang2023math,wizardmath}. Here, we present Group Relative Policy Optimization (GRPO), a novel RL approach that is both efficient and effective.

\subsubsection{From PPO to GRPO} Proximal Policy Optimization (PPO) \citep{schulman2017proximal} is an actor-critic RL algorithm that is widely used in the RL fine-tuning stage of LLMs \citep{ouyang2022training}. In particular, it optimizes LLMs by maximizing the following surrogate objective:

\begin{equation}
\footnotesize
    \mathcal{J}_{PPO}(\theta) = \mathbb{E}{[q \sim P(Q), o \sim \pi_{\theta_{old}}(O|q)]} \frac{1}{|o|} \sum_{t=1}^{|o|} \min \left[ \frac{\pi_\theta(o_{t} | q, o_{<t})}{\pi_{\theta_{old}}(o_{t} | q, o_{<t})} A_{t}, \text{clip} \left( \frac{\pi_\theta(o_{t} | q, o_{<t})}{\pi_{\theta_{old}}(o_{t} | q, o_{<t})}, 1 - \varepsilon, 1 + \varepsilon \right)  A_{t} \right] ,
\end{equation}

The objective function utilized in PPO, $\mathcal{J}_{PPO}(\theta)$, is formalized as follows. We consider the expectation over queries $q$ sampled from the distribution $P(Q)$ and output sequences $o$ generated by the previous policy $\pi_{\theta_{old}}(O|q)$. For each token in the output, indexed by $t$ across the output length $|o|$, the objective aggregates the minimum between the scaled advantage using the current and old policy ratio, and the advantage multiplied by the same ratio that is restricted to the interval $[1-\varepsilon, 1+\varepsilon]$. This clipping mechanism is designed to penalize deviations of the new policy from the old one beyond a preset threshold, promoting stability during updates.

Here, $\pi_{\theta}$ denotes the updated policy model while $\pi_{\theta_{old}}$ refers to the previous policy version. The variables $q$ and $o$ represent questions sampled from the dataset and corresponding outputs produced by $\pi_{\theta_{old}}$. The hyper-parameter $\varepsilon$ is employed for clipping in Proximal Policy Optimization (PPO) to promote stable optimization. The term $A_t$ stands for the advantage function, which is estimated using Generalized Advantage Estimation (GAE) \citep{gae}, relying on the rewards $\{r_{\ge t}\}$ and a parameterized value function $V_{\psi}$. Consequently, PPO requires concurrent training of the value function alongside the policy model. To prevent excessive reward model optimization, a common strategy is to incorporate a KL divergence penalty between the policy and a reference model at every token into the reward, as proposed by \citep{ouyang2022training}; that is,

\begin{equation}
    r_{t} = r_\varphi(q, o_{\le t}) - \beta \log\frac{\pi_{\theta}(o_{t}|q, o_{<t})}{\pi_{ref}(o_{t}|q, o_{<t})},
\label{eq:PPO-reward}
\end{equation}

Here, $r_\varphi$ denotes the reward model, $\pi_{ref}$ refers to the reference model—typically the initial SFT model—and $\beta$ represents the weight assigned to the KL divergence penalty.

As the value function employed in PPO is typically another model of comparable size as the policy model, it brings a substantial memory and computational burden. Additionally, during RL training, the value function is treated as a baseline in the calculation of the advantage for variance reduction. While in the LLM context, usually only the last token is assigned a reward score by the reward model, which may complicate the training of a value function that is accurate at each token. To address this, as shown in Figure \ref{fig:grpo}, we propose Group Relative Policy Optimization (GRPO), which obviates the need for additional value function approximation as in PPO, and instead uses the average reward of multiple sampled outputs, produced in response to the same question, as the baseline. More specifically, for each question $q$, GRPO samples a group of outputs $\{o_1, o_2, \cdots, o_G\}$  from the old policy  $\pi_{\theta_{old}}$  and then optimizes the policy model by maximizing the following objective:

\begin{equation}
\footnotesize
\begin{split}
    \mathcal{J}_{GRPO}(\theta) &= \mathbb{E}{[q \sim P(Q), \{o_i\}_{i=1}^G \sim \pi_{\theta_{old}}(O|q)]}  \\
    & \frac{1}{G}\sum_{i=1}^G\frac{1}{|o_i|} \sum_{t=1}^{|o_i|} \Bigg\{ \min \Bigg[ \frac{\pi_\theta(o_{i,t} | q, o_{i,<t})}{\pi_{\theta_{old}}(o_{i,t} | q, o_{i,<t})} \hat{A}_{i,t},~ \text{clip} \left( \frac{\pi_\theta(o_{i,t} | q, o_{i,<t})}{\pi_{\theta_{old}}(o_{i,t} | q, o_{i,<t})},~ 1 - \varepsilon,~ 1 + \varepsilon \right)  \hat{A}_{i,t} \Bigg] \\
    & \qquad - \beta~ \mathbb{D}_{KL}\left[\pi_{\theta} || \pi_{ref}\right] \Bigg\},
\end{split}
\label{eq:GRPO-obj}
\end{equation}

where $\varepsilon$ and $\beta$ are hyper-parameters, and $\hat{A}_{i,t}$  is the advantage calculated based on relative rewards of the outputs inside each group only, which will be detailed in the following subsections. The group relative way that GRPO leverages to calculate the advantages, aligns well with the comparative nature of rewards models,  as reward models are typically trained on datasets of comparisons between outputs on the same question. Also note that, instead of adding KL penalty in the reward, GRPO regularizes by directly adding the KL divergence between the trained policy and the reference policy to the loss, avoiding complicating the calculation of  $\hat{A}_{i,t}$. And different from the KL penalty term used in (\ref{eq:PPO-reward}), we estimate the KL divergence with the following unbiased estimator \citep{kl_approx}:

\begin{equation}
\small
    \mathbb{D}_{KL}\left[\pi_{\theta} || \pi_{ref}\right] = \frac{\pi_{ref}(o_{i,t}|q,o_{i,<t})}{\pi_{\theta}(o_{i,t}|q,o_{i,<t})} - \log\left(\frac{\pi_{ref}(o_{i,t}|q,o_{i,<t})}{\pi_{\theta}(o_{i,t}|q,o_{i,<t})}\right) - 1,
\end{equation}

which is guaranteed to be positive.

\begin{algorithm}[t]
  \small
  \caption{Iterative Group Relative Policy Optimization}
  \textbf{Input} initial policy model $\pi_{\theta_{\text{init}}}$; reward models $r_\varphi$; task prompts $\mathcal{D}$; 
  hyperparameters $\varepsilon$, $\beta$, $\mu$
  \begin{algorithmic}[1]
    \State Initialize policy model: $\pi_\theta \leftarrow \pi_{\theta_{\text{init}}}$
    \For{iteration = 1, \dots, I}
      \State Set reference model: $\pi_{ref} \leftarrow \pi_{\theta}$
      \For{step = 1, \dots, M}
      \State Draw a minibatch $\mathcal{D}_b$ from $\mathcal{D}$
      \State Assign old policy: $\pi_{\theta_{old}} \leftarrow \pi_{\theta}$
      \State For each query $q$ in $\mathcal{D}_b$, generate $G$ responses $\{o_i\}_{i=1}^G \sim \pi_{\theta_{old}} (\cdot \mid q)$
      \State Evaluate rewards $\{r_i\}_{i=1}^{G}$ corresponding to each response $o_i$ using $r_{\varphi}$
      \State Estimate group-relative advantages $\hat{A}_{i,t}$ for the $t$-th token of each output $o_i$
      \For{GRPO iteration = 1, \dots, $\mu$}
        \State Optimize the policy $\pi_{\theta}$ with respect to the GRPO objective (Equation \ref{eq:GC-GRPO})
      \EndFor
    \EndFor
    \State Continue training $r_\varphi$ with replay buffer updates.
    \EndFor
  \end{algorithmic}
  \textbf{Output} $\pi_\theta$
  \label{alg:iter-grpo}
\end{algorithm}

\subsubsection{Outcome Supervision RL with GRPO} In a formal sense, given a question $q$, the approach involves generating a batch of outputs $\{o_1, o_2, \cdots, o_G\}$ using the previous iteration of the policy model $\pi_{\theta_{old}}$. Each of these candidate outputs is assessed by a reward model, which assigns a set of $G$ reward values $\mathbf{r}=\{r_1, r_2, \cdots, r_G\}$. These obtained rewards are standardized by first removing the mean of the group and then dividing by the standard deviation, resulting in normalized scores. Within outcome supervision, each final normalized reward is provided at the conclusion of the corresponding output $o_i$, and this value is used as the advantage $\hat{A}_{i, t}$ for every token within the sequence, i.e., $\hat{A}_{i, t} = \widetilde{r}_i = \frac{r_i- {\rm mean}(\mathbf{r})}{{\rm std}(\mathbf{r})}$. The policy model is subsequently refined by maximizing the objective specified in equation (\ref{eq:GRPO-obj}).

\subsubsection{Process Supervision RL with GRPO} Unlike outcome supervision, which offers a single reward signal upon generating an output, process supervision supplies feedback at each reasoning step, leading to more granular and potentially more effective guidance for solving complex mathematical problems. As in \cite{wang2023math}, we investigate this process-based supervisory approach. Specifically, for any given question $q$ and a set of $G$ sampled solutions $\{o_1, o_2, \cdots, o_G\}$, a process reward model assigns a score to every individual reasoning step within each solution. This produces the reward set $\mathbf{R} = \{ \{r_1^{index(1)},\cdots,r_1^{index(K_1)}\}, \cdots,  \{r_G^{index(1)},\cdots,r_G^{index(K_G)}\} \}$, where $index(j)$ refers to the token index terminating the $j$-th reasoning step, and $K_i$ indicates how many steps comprise the $i$-th solution. We then normalize these rewards by subtracting the mean and dividing by the standard deviation across all rewards: $\widetilde{r}_i^{index(j)} = \frac{r_i^{index(j)} - {\rm mean(\mathbf{R})}}{{\rm std(\mathbf{R})}}$.
Next, process supervision estimates each token’s advantage as the cumulative normalized reward from all subsequent reasoning steps, formalized as $\hat{A}_{i, t} = \sum_{index(j) \ge t} \widetilde{r}_i^{index(j)}$.
The policy is then improved by maximizing the objective function set forth in equation (\ref{eq:GRPO-obj}).

\subsubsection{Iterative RL with GRPO} During the reinforcement learning training, the previously learned reward model may become inadequate for guiding the evolving policy model. To address this, we investigate the use of iterative RL with GRPO. As detailed in Algorithm \ref{alg:iter-grpo}, the iterative GRPO approach involves creating updated training datasets for the reward model by leveraging outputs sampled from the current policy model, and continuously updating the reward model through a replay strategy that retains 10\% of prior data. Subsequently, we designate the policy model as the reference model and proceed to iteratively refine the policy model with guidance from the updated reward model.

\subsection{Training and Evaluating DeepSeekMath-RL}

We implement reinforcement learning (RL) using DeepSeekMath-Instruct 7B as the base model. The RL training dataset is comprised of approximately 144K chain-of-thought-style questions associated with GSM8K and MATH, drawn from the SFT data. To specifically assess the effects of RL on benchmarks lacking corresponding data during RL, we omit other SFT questions from the RL phase. The reward model’s training data is prepared according to the protocol outlined in \citep{wang2023math}. The reward model is initially trained on DeepSeekMath-Base 7B, employing a learning rate of 2e-5. The policy model for GRPO is optimized with a learning rate of 1e-6 and a KL divergence coefficient of 0.04. For each individual question, we generate $64$ responses. Both the maximum sequence length and the training batch size are fixed at 1024. Following every exploration phase, the policy model undergoes a single update. For evaluation, DeepSeekMath-RL 7B is benchmarked following the same suite as DeepSeekMath-Instruct 7B. Within DeepSeekMath-RL 7B, GSM8K and MATH—both employing chain-of-thought reasoning—are classified as in-domain, while all remaining benchmarks are considered out-of-domain tasks.

Table \ref{tab:sft_rl_math} presents results comparing open- and closed-source models employing both chain-of-thought and tool-augmented reasoning approaches across English and Chinese datasets. Key observations include: 1) DeepSeekMath-RL 7B achieves accuracy rates of 88.2\% on GSM8K and 51.7\% on MATH using chain-of-thought reasoning. These results outperform all open-source models within the 7B to 70B parameter range, as well as most closed-source counterparts. 2) Importantly, DeepSeekMath-RL 7B’s training is limited to instruction tuning with chain-of-thought-formatted data from GSM8K and MATH, initialized from DeepSeekMath-Instruct 7B. Even within this narrowly defined training set, it consistently exceeds the performance of DeepSeekMath-Instruct 7B across all measured criteria, highlighting the notable advantages of reinforcement learning.

\section{Discussion}
This section presents an overview of the results obtained from both our pre-training phase and reinforcement learning (RL) experiments.

\subsection{Lessons Learnt in Pre-Training}

To begin, we describe our pre-training process. Except where noted, our methodology follows the configuration presented in Section \ref{sec:quality-policy}. Additionally, throughout this section, mentions of the DeepSeekMath Corpus pertain to an 89B-token dataset derived from the second round of data collection.

\subsubsection{Code Training Benefits Mathematical Reasoning}

A widely held, though not empirically confirmed, conjecture posits that training on code enhances reasoning skills. In this work, we present an initial investigation into this claim, focusing specifically on mathematics: training with code bolsters models' mathematical reasoning capacities, regardless of whether they utilize external tools.

To study how code training affects mathematical reasoning, we experimented with the following two-stage training and one-stage training settings:

\textbf{Two-Stage Training}

\begin{itemize}[topsep=0pt]
    \item \textbf{Code Training for 400B Tokens $\rightarrow$ Math Training for 150B Tokens}: Our approach involves initially pretraining DeepSeek-LLM 1.3B on a dataset of 400B code tokens, and subsequently fine-tuning it with 150B math-specific tokens;
    \item \textbf{General Training for 400B Tokens $\rightarrow$ Math Training for 150B Tokens}: To serve as a baseline, we conduct an additional experiment in which general tokens—sampled from a comprehensive general-purpose corpus assembled by DeepSeek-AI—replace code tokens during the first training stage. This setup is designed to explore whether code data offers greater benefits than general data for enhancing mathematical reasoning abilities.
\end{itemize}

\textbf{One-Stage Training}

\begin{itemize}[topsep=0pt]
    \item \textbf{Math Training Utilizing 150B Tokens}: DeepSeek-LLM 1.3B is trained exclusively on 150B math tokens;
    \item \textbf{Joint Training with 400B Code Tokens and 150B Math Tokens}: Sequentially training on math data after code data results in diminished code capabilities. To address this, we explore if incorporating both code and math tokens concurrently in a single training phase can enhance mathematical reasoning while mitigating catastrophic forgetting in code-related skills.
\end{itemize}

\paragraph{Results}
Tables \ref{tab:code-math} and \ref{tab:code-math-general-eval} present the downstream outcomes across various training configurations.

Code training substantially enhances program-aided mathematical reasoning, applicable to both two-stage and one-stage training paradigms. As illustrated in Table \ref{tab:code-math}, within the two-stage framework, applying code training independently already leads to marked improvements in solving GSM8K and MATH problems via Python. Subsequent math training in the latter phase further augments these gains. Notably, under the one-stage training paradigm, incorporating both code and math tokens concurrently effectively addresses the problem of catastrophic forgetting seen in two-stage training, while also producing complementary benefits for both coding ability (refer to Table \ref{tab:code-math-general-eval}) and program-aided mathematical reasoning (refer to Table \ref{tab:code-math}).

Training on code additionally enhances mathematical reasoning even in the absence of tool assistance. In the two-stage training framework, the preliminary phase involving code training yields notable, though moderate, improvements. This phase also facilitates more effective math training in the next stage, culminating in optimal performance outcomes. Conversely, adopting a single-stage training process where both code and math tokens are mixed appears to detract from mathematical reasoning when tool use is not available. One possible explanation is that DeepSeek-LLM 1.3B may not possess sufficient capacity, given its relatively small size, to thoroughly integrate both code and mathematical data at the same time.

\subsubsection{ArXiv Papers Seem Ineffective in Improving Mathematical Reasoning}

ArXiv papers are commonly included as a component of math pre-training data \citep{minerva,gpt-f,llemma,mathpile}. However, detailed analysis regarding their impact on mathematical reasoning has not been extensively conducted. Perhaps counter-intuitively, according to our experiments, arXiv papers seem ineffective in improving mathematical reasoning. We experiment with models of different sizes, including DeepSeek-LLM 1.3B and DeepSeek-Coder-Base-v1.5 7B \citep{deepseek-coder}, using arXiv corpora that underwent varied processing pipelines:

\begin{itemize}[topsep=0pt]
    \item \textbf{MathPile} \citep{mathpile}: a dataset containing 8.9B tokens, curated through heuristic-based cleaning and filtering, with scientific arXiv articles making up over 85\% of its content;
    \item \textbf{ArXiv-RedPajama} \citep{redpajama}: a 28.0B-token collection comprising all arXiv LaTeX source files, systematically stripped of preambles, comments, macros, and bibliographic entries.
\end{itemize}

During our experiments, we trained DeepSeek-LLM 1.3B on 150B tokens and DeepSeek-Coder-Base-v1.5 7B on 40B tokens using each separate arXiv corpus. Results indicate that utilizing arXiv papers alone does not enhance mathematical reasoning abilities. When models are exposed exclusively to the arXiv-only dataset, there are either minimal gains or observable declines in performance on a range of mathematical evaluation tasks differing in complexity within this study. The evaluated benchmarks encompass quantitative reasoning sets such as GSM8K and MATH (refer to Table \ref{tab:arxiv-cot}), multiple-choice tests exemplified by MMLU-STEM (Table \ref{tab:arxiv-cot}), and formal mathematics problems like those in miniF2F (see Table \ref{tab:arxiv_atp}).

However, this conclusion has its limitations and should be taken with a grain of salt. We have not yet studied:

\begin{itemize}[topsep=0pt]
    \item How arXiv tokens might influence additional mathematics-focused tasks beyond the scope of this study, for example, the transformation of formal theorems or proofs into their informal counterparts;
    \item The interaction between arXiv tokens and alternative data sources and its resulting impact;
    \item If the observed advantages of incorporating arXiv papers would persist or amplify in the context of larger-scale models.
\end{itemize}

Thus, further exploration is required, which we leave for future studies.

\subsection{Insights of Reinforcement Learning}
\subsubsection{Towards to a Unified Paradigm} In this section, we introduce a cohesive framework to examine various training strategies, including SFT, RFT, DPO, PPO, and GRPO, and perform experiments to investigate the key components of this unified approach. Typically, for any training method, the gradient with respect to the model parameters $\theta$ can be expressed as:

\begin{equation}
    \nabla_{\theta}\mathcal{J}_{\textcolor{red}{\mathcal{A}}}(\theta) = \mathbb{E}[\underbrace{(q,o) \sim \textcolor{red}{\mathcal{D}}}_{Data \ Source}]\left( \frac{1}{|o|} \sum_{t=1}^{|o|}  \underbrace{GC_{{\mathcal{A}}}(q, o, t, \textcolor{red}{\pi_{{rf}}})}_{Gradient \ Coefficient}  \nabla_{\theta}\log \pi_{\theta}(o_t | q, o_{<t})\right).
\label{eq:objective}
\end{equation}

There exist three key components: 1) \textit{Data Source $\mathcal{D}$}, which determines the training data; 2) \textit{Reward Function $\pi_{{rf}}$}, which is the source of the training reward signal; 3) \textit{Algorithm $\mathcal{A}$}: which processes the training data and the reward signal to the gradient coefficient $GC$ that determines the magnitude of the penalty or reinforcement for the data. We analyze several representative methods based on such a unified paradigm:

\begin{itemize}
    \item  \textbf{Supervised Fine-tuning (SFT)}: SFT involves adapting a pretrained model to specific tasks by training on SFT datasets curated by human annotators.
    \item \textbf{Rejection Sampling Fine-tuning (RFT)}: RFT extends SFT by utilizing responses generated by the SFT model for SFT prompts, and subsequently retaining only those responses deemed correct before further fine-tuning.
    \item \textbf{Direct Preference Optimization (DPO)}: DPO enhances the SFT model by training on additional outputs generated from the SFT model itself, applying a pair-wise DPO loss for further optimization.
    \item \textbf{Online Rejection Sampling Fine-tuning (Online RFT)}: Unlike standard RFT, Online RFT continually updates the policy model, beginning from the SFT model, and improves it by leveraging augmented outputs dynamically sampled from the ongoing policy model.
    \item \textbf{PPO/GRPO}: In this approach, the SFT model serves as the initial policy, which is subsequently strengthened via outputs interactively generated by the continually updated policy model.
\end{itemize}

A summary of the constituent elements of these approaches is presented in Table \ref{tab:data_policy}. For a comprehensive explanation of the derivation steps, please consult Appendix \ref{app:analysis-rl}.

\paragraph{Observation about Data Source} The data source is classified into two types: online sampling and offline sampling. In online sampling, training data is collected from the outcomes generated by the policy model during its current training process. In contrast, offline sampling refers to using training data derived from the results of the original SFT model. Methods such as RFT and DPO utilize the offline approach, whereas Online RFT and GRPO adopt an online approach.

As illustrated in Figure \ref{fig:rl-analysis}, our results indicate that Online RFT achieves markedly better performance than RFT on both benchmarks. In particular, during the early training phase, Online RFT performs similarly to RFT; however, as training progresses, Online RFT establishes a clear lead, highlighting the effectiveness of online training. This observation aligns with expectations: early in training, the behavior of the actor is still closely aligned with that of the SFT model, resulting in sampled data with minimal discrepancies. As training advances, the data drawn from the actor tends to diverge more substantially, at which point the benefits of on-the-fly data sampling become more pronounced.

\paragraph{Observation about Gradient Coefficient} Within the algorithm, the reward signal is transformed into a gradient coefficient that directs the adjustment of model parameters. In our experimental setup, the reward function is bifurcated into `Rule' and `Model' components. The `Rule' aspect evaluates response quality through the lens of answer correctness, whereas the `Model' aspect involves developing a reward model that assigns a score to each response; the training data for this model derives from the rule-based assessments. As shown in Equations \ref{eq:GC-RFT} and \ref{eq:GC-GRPO}, there is a notable distinction between GRPO and Online RFT. Specifically, GRPO is characterized by its dynamic modification of the gradient coefficient in accordance with the score yielded by the reward model, facilitating tailored reinforcement or punishment contingent on the response's score. In opposition, Online RFT does not share this capability: it refrains from imposing penalties for incorrect responses and consistently applies the same level of reinforcement to all responses deemed correct, regardless of their relative quality.

As illustrated in Figure \ref{fig:rl-analysis}, GRPO outperforms the online RFT, underscoring the effectiveness of adjusting positive and negative gradient coefficients. Moreover, GRPO+PS achieves better results than GRPO+OS, which demonstrates the advantage of employing step-specific, fine-grained gradient coefficients. Additionally, we investigate the impact of iterative RL; in our setup, we apply two iterations. Figure \ref{fig:iter-rl} reveals that iterative RL leads to notable performance gains, with the most substantial improvement occurring during the initial iteration.

\subsubsection{Why RL Works?} In this study, reinforcement learning is applied using a portion of instruction tuning datasets, which leads to marked improvements over models solely utilizing instruction tuning. To further clarify the mechanism behind the effectiveness of reinforcement learning, we assess Pass@K and Maj@K accuracy metrics for both Instruct and RL models across two different evaluation benchmarks. Figure \ref{fig:maj-pass} illustrates that while RL substantially increases Maj@K accuracy, it does not impact Pass@K scores. These results suggest that RL contributes to making the output distribution more resilient, implying that \textbf{the observed gains arise primarily from an increased likelihood of the correct answer being present among the TopK responses rather than from advances in intrinsic model capability.} In a related context, \citep{wang2023making} describe a \textbf{misalignment problem} in reasoning-related tasks for SFT models, and further demonstrate that adopting various preference alignment techniques can enhance the reasoning performance of these models \citep{yuan2023rrhf,song2023preference,wang2023making}.

\subsubsection{How to Achieve More Effective RL?}
\label{sec:effec_rl}
Our findings indicate that RL performs robustly in the domain of mathematical reasoning. Additionally, we propose a cohesive framework that integrates a range of prominent training approaches. According to this framework, all examined methods can be interpreted as variations or streamlined forms of RL strategies. As outlined in Equation \ref{eq:objective}, three primary elements are identified: Data Source, Algorithm, and Reward Function. We discuss prospective avenues for improvement regarding these three aspects.

\paragraph{Data Source} The data source constitutes the foundational input for all training frameworks. Within RL scenarios, we define the data source as the set of unlabeled questions, accompanied by outputs generated from the policy model. In this study, our experiments are limited to questions drawn exclusively from the instruction tuning phase, with responses obtained through basic nucleus sampling. We suspect that this methodological constraint might contribute to our RL pipeline’s improvements being restricted to Maj@K metrics. Going forward, we plan to investigate our RL pipeline using out-of-distribution question prompts and to incorporate \textbf{advanced sampling (decoding) strategies}—such as tree-search-based approaches \citep{yao2023tree}. Furthermore, \textbf{efficient inference techniques} \citep{xia-etal-2023-speculative,leviathan2023fast,kwon2023efficient,xia2024unlocking}, which govern how effectively the policy model can explore, are anticipated to be of significant importance as well.

\paragraph{Algorithms} Algorithms utilize both data and reward signals to determine the gradient coefficients, which in turn guide the update of model parameters. According to Equation \ref{eq:objective}, current approaches essentially place full \textbf{TRUST} in the reward function's signal, using it to modulate the conditional probability of specific tokens. Nonetheless, the reliability of the reward signal cannot be guaranteed at all times, particularly when tackling tasks with substantial complexity. For instance, the PRM800K datasets \citep{lightman2023let}—despite being meticulously curated by expert annotators—still exhibit a misannotation rate of around 20\%\footnote{\url{https://github.com/openai/prm800k/issues/12\#issuecomment-1728491852}}. Therefore, we aim to investigate reinforcement learning algorithms designed to be resilient to noisy reward feedback. We contend that \textbf{WEAK-TO-STRONG} alignment methods \citep{burns2023weak} represent a potentially transformative advancement in the underlying learning algorithms.

\paragraph{Reward Function} The reward function serves as the principal mechanism for providing training signals. In RL, it is commonly instantiated as a neural reward model. We identify three key areas of advancement for reward models: 1) \textbf{Strategies to improve the generalization capacity of the reward model.} A robust reward model should generalize well to unseen (out-of-distribution) queries and more sophisticated decoding outcomes; failing to achieve this may cause reinforcement learning to simply stabilize LLM distribution without genuinely advancing core performance; 2) \textbf{Methods to represent the uncertainty inherent in the reward model.} Capturing uncertainty can create vital connections between imperfect reward models and weak-to-strong learning frameworks; 3) \textbf{Approaches to construct high-quality, process-level reward models efficiently} so that they can deliver nuanced training feedback throughout the reasoning chain \citep{lightman2023let,wang2023math}.

\section{Conclusion, Limitation, and Future Work}

In this work, we introduce DeepSeekMath, a model that surpasses all existing open-source systems on the MATH benchmark for competition-level problems, and demonstrates performance on par with that of proprietary models. DeepSeekMath builds upon DeepSeek-Coder-v1.5 7B as its base and is continually trained over 500B tokens, of which 120B are mathematics-focused tokens primarily drawn from Common Crawl. Our thorough ablation analysis reveals that web pages represent a particularly fruitful resource for high-quality mathematical content, whereas data from arXiv does not deliver the anticipated benefits. We propose Group Relative Policy Optimization (GRPO), a modification of Proximal Policy Optimization (PPO), which achieves substantial improvements in mathematical reasoning, while also reducing memory requirements. Experimental findings indicate that GRPO remains advantageous, even when DeepSeekMath-Instruct 7B attains strong performance on evaluation benchmarks. Furthermore, we propose a unified framework to interpret various methods, and outline several promising pathways for advancing reinforcement learning techniques.

While DeepSeekMath~demonstrates strong results on benchmarks involving quantitative reasoning, its performance lags behind that of closed models when it comes to geometry and theorem-proving tasks. For example, during preliminary testing, the model failed to solve questions involving triangles and ellipses, which suggests a possible bias in the data utilized during pre-training and fine-tuning stages. Furthermore, due to its scale, DeepSeekMath~underperforms compared to GPT-4 in few-shot scenarios; while GPT-4 benefits from few-shot examples to enhance accuracy, DeepSeekMath~exhibits little difference between zero-shot and few-shot settings. Going forward, we intend to enhance our engineered data selection process to assemble a more comprehensive and higher quality pre-training dataset. Additionally, we plan to investigate avenues outlined in Section \ref{sec:effec_rl} for reinforcing learning approaches that yield more effective LLM training.

\end{document}